
\begin{frame}{Feedback Loops: The Model Changes Its Own Training Data}

    \centering
    \begin{tikzpicture}[
        node distance=1.7cm and 2.4cm,
        box/.style={rectangle, draw, rounded corners, minimum height=1.0cm, text width=3.0cm,
                    align=center, font=\small},
        arr/.style={-{Latex[length=2.5mm]}, thick}
    ]
        \node[box, fill=raiBlue!12] (model)  {Model predicts high risk in area A};
        \node[box, fill=gray!12, right=of model] (action) {More patrols sent to area A};
        \node[box, fill=raiRed!12, below=of action] (record) {More arrests recorded in area A};
        \node[box, fill=raiRed!20, below=of model]  (retrain) {Retraining confirms area A is high risk};

        \draw[arr] (model)  -- (action);
        \draw[arr] (action) -- (record);
        \draw[arr] (record) -- (retrain);
        \draw[arr] (retrain) -- (model);
    \end{tikzpicture}

    \vspace{0.8em}

    \begin{itemize}\small
        \setlength\itemsep{0.35em}
        \item The loop is self-confirming: the prediction generates the evidence for itself.
        \item Same shape in recommenders (you only observe clicks on what you showed),
              in credit (you only observe repayment for approved applicants), and in hiring.
        \item \textbf{Detection:} hold out a slice from the policy, or log what the model
              \emph{suppressed}, not just what it surfaced.
    \end{itemize}
\end{frame}
